\subsection{PRINCIPAL PART OF A PRIMITIVE}

Let \(f\) be a regulated real function \(\neq 0\) having constant sign on an interval \([a, +\infty[\) ; the following proposition allows one to obtain the principal part of the primitive \(\int_{x}^{+\infty} f(t) dt\) if \(\int_{a}^{+\infty} f(t) dt\) converges in certain cases, and of the primitive \(\int_{a}^{x} f(t) dt\) if \(\int_{a}^{+\infty} f(t) dt\) is infinite:

PROPOSITION 8. Put \(\mathrm{F}(x) = \int_{x}^{+\infty} f(t) dt\) if \(\int_{a}^{+\infty} f(t) dt\) converges, and \(\mathrm{F}(x) = \int_{a}^{x} f(t) dt\) if \(\int_{a}^{+\infty} f(t) dt\) is infinite. Suppose that \(\log |f|\) and \(\log x\) are comparable of order 1.

\(1^{\circ}\) If \(f\) is of finite order \(\mu \neq -1\) with respect to \(x\) one has

\[
\mathrm{F} (x) \sim \frac {1}{| \mu + 1 |} x f (x).\tag{1}
\]

\(2^{\circ}\) If \(f\) is of infinite order relative to \(x\) and if \(f / f'\) and \(x\) are comparable of order 1, then one has

\[
\mathrm{F} (x) \sim \frac {(f (x)) ^ {2}}{| f ^ {\prime} (x) |}.\tag{2}
\]

One should note that the hypothesis implies that \(f(x)\) has a determinate order relative to \(x\) (V, p.~219).

\(1^{\circ}\) If \(f\) is of order \(\mu \neq 0\) relative to \(x\) one has \(\log |f| \sim \mu \log x\) , so, since \(\log |f|\) and \(\log x\) are comparable of order 1, by prop. 7 of V, p.~232 one has \(f'/f \sim\) \(\mu/x\) , whence \(xf' \sim \mu f\) . If \(\mu > -1\) one has \(f(x) \gg x^{\mu-\varepsilon}\) for every \(\varepsilon > 0\) , and hence (V, p.~229, prop. 2) the integral \(\int_{a}^{+\infty} f(t) dt\) is infinite. One can write \(F(x) = \int_{a}^{x} f(t) dt = xf(x) - af(a) - \int_{a}^{x} tf'(t) dt\) , whence again

\[
\int_ {a} ^ {x} \left(f (t) + t f ^ {\prime} (t)\right) d t = x f (x) - a f (a);
\]

since \(\mu \neq -1\) we have \(f(x) + xf'(x) \sim (\mu + 1)f(x)\) , so (V, p.~230, prop. 6)

\[
\int_ {a} ^ {x} \left(f (t) + t f ^ {\prime} (t)\right) d t \sim (\mu + 1) \mathrm{F} (x),
\]

which proves (1) in this case. If \(\mu = 0\) one has likewise that \(xf'(x) \ll f(x)\) , which again gives \(f(x) + xf'(x) \sim f(x)\) . One argues similarly when \(\mu < -1\) , in the case where \(\int_{a}^{+\infty} f(t) dt\) converges.

\(2^{\circ}\) If \(f\) is of order \(+\infty\) relative to \(x\) one has \(\log |f| \gg \log x\) , so (V, p.~232, prop. 7) \(f'/f \gg 1/x\) , or again, putting \(g(x) = f(x)/f'(x)\) , \(g(x) \ll x\) ; further, since \(f(x) \gg x^{\alpha}\) for every \(\alpha > 0\) the integral \(\int_{a}^{+\infty} f(t) dt\) is infinite. One can write

\[
\mathrm{F} (x) = \int_ {a} ^ {x} f (t) d t = \int_ {a} ^ {x} g (t) f ^ {\prime} (t) d t = g (x) f (x) - g (a) f (a) - \int_ {a} ^ {x} f (t) g ^ {\prime} (t) d t;
\]

since \(g\) and \(x\) are comparable of order 1, from the relation \(g(x) \ll x\) one can deduce (V, p.~232, prop. 7) that \(g'(x) \ll 1\) , hence that \(fg' \ll f\) , and consequently (V, p.~230, prop.6)

\[
\int_ {a} ^ {x} f (t) g ^ {\prime} (t) d t \ll \mathrm{F} (x),
\]

which establishes the relation (2). The proof is similar when \(f\) is of order \(-\infty\) relative to \(x\) , in the case where \(\int_{a}^{+\infty} f(t) dt\) converges.

Let E be a comparison scale (for real x tending to \(+\infty\) ) formed of nonzero real functions that are of constant sign on a neighbourhood of \(+\infty\) , such that \(x \in E\) and such that the product and quotient of two functions in E again belongs to E (V, p.~221 and p.~224). If a regulated function f with constant sign on a neighbourhood of \(+\infty\) has a principal part cg relative to E, then \(\int_{x}^{+\infty} f(t) dt\) (resp. \(\int_{a}^{x} f(t) dt\) according to the case) will be equivalent to \(c \int_{x}^{+\infty} g(t) dt\) (resp. \(c \int_{a}^{x} g(t) dt\) ); if the function g satisfies the hypotheses of prop. 8 of V, p.~233, and if (when the formula (2) of V, p.~233, applies) one knows a principal part of \(g'\) relative to E, one will thus have a principal part of \(\int_{x}^{+\infty} f(t) dt\) (resp. \(\int_{a}^{x} f(t) dt\) ) relative to E.

Examples. 1) The function \(1 / \log x\) is of order 0 relative to \(x\) and satisfies the conditions of prop. 8 of V, p.~233; thus

\[
\int_ {a} ^ {x} \frac {d t}{\log t} \sim \frac {x}{\log x}.
\]

\begin{enumerate}
\def\labelenumi{\arabic{enumi})}
\setcounter{enumi}{1}
\tightlist
\item
  The function \(e^{x^2}\) is of order \(+\infty\) relative to \(x\) and satisfies the conditions of prop. 8, so
\end{enumerate}

\[
\int_ {a} ^ {x} e ^ {t ^ {2}} d t \sim \frac {1}{2 x} e ^ {x ^ {2}}.
\]

In the Appendix (V, p.~252) we shall define a set of functions for which prop. 7 and prop. 8 always apply.

Remark. Prop. 8 does not apply directly to a function \(f\) of order \(-1\) relative to \(x\) . But then one can write \(f(x) = f_1(x) / x\) , with \(f_1\) of order 0 relative to \(x\) . Suppose, for example, that \(\int_{a}^{+\infty} f(t) dt\) is infinite; then

\[
\mathrm{F} (x) = \int_ {a} ^ {x} f (t) d t = \int_ {a} ^ {x} \frac {1}{t} f _ {1} (t) d t = \int_ {\log a} ^ {\log x} f _ {1} \left(e ^ {u}\right) d u.
\]

If the function \(f_{1}(e^{y})\) satisfies the hypotheses of prop. 8 and has an order \(\neq -1\) relative to \(y\) (that is, if \(f_{1}(x)\) has an order \(\neq -1\) relative to \(\log x\) ) the formulae (1) and (2) again allow us to obtain a principal part for \(F(x)\) . For example, let \(f(x) = \frac{\exp(\sqrt{\log x})}{x\log x}\) ; since \(\exp(\sqrt{\log x})\) is of order 0 relative to \(x\) , \(f\) is of order \(-1\) ; here one has \(f_{1}(e^{y}) = e^{\sqrt{y}} / y\) , and this function is of order \(+\infty\) relative to \(y\) ; prop. 8 applies and gives \(\int_{\alpha}^{y} e^{\sqrt{u}} / u du \sim 2e^{\sqrt{y}} / \sqrt{y}\) ; on reverting to the variable \(x\) we have \(\int_{a}^{x} \frac{\exp(\sqrt{\log t})}{t\log t} dt \sim \frac{2\exp(\sqrt{\log x})}{\sqrt{\log x}}\) .

